\documentclass[11pt,a4paper]{article}
\usepackage[margin=25mm,headheight=15pt]{geometry}
\usepackage{fontspec}
\setmainfont{Latin Modern Roman}
\setsansfont{Latin Modern Sans}
\setmonofont{DejaVu Sans Mono}[Scale=0.80]
\usepackage{amsmath,amssymb,amsthm,mathtools}
\usepackage{microtype,booktabs,longtable,array,tabularx}
\usepackage[dvipsnames]{xcolor}
\definecolor{ForgeBlue}{HTML}{17344C}
\definecolor{ForgeGray}{HTML}{566573}
\definecolor{ForgePale}{HTML}{F3F6F8}
\usepackage{enumitem,fvextra,tcolorbox}
\tcbuselibrary{breakable}
\usepackage{fancyhdr,lastpage,titlesec}
\usepackage{hyperref,xurl}
\hypersetup{colorlinks=true,linkcolor=ForgeBlue,citecolor=ForgeBlue,urlcolor=ForgeBlue,
 pdftitle={Forge: A Certificate-First Design Beyond Lean's grind},
 pdfauthor={OpenAI},pdfsubject={Lean tactic architecture and tested Python research prototypes}}
\pagestyle{fancy}
\fancyhf{}\fancyhead[L]{\sffamily\small FORGE / PROOF AUTOMATION}
\fancyhead[R]{\sffamily\small Design and prototype report}
\fancyfoot[L]{\sffamily\footnotesize 14 September 2026}
\fancyfoot[R]{\sffamily\footnotesize \thepage\ / \pageref*{LastPage}}
\renewcommand{\headrulewidth}{0.3pt}
\titleformat{\section}{\Large\sffamily\bfseries\color{ForgeBlue}}{\thesection}{0.65em}{}
\titleformat{\subsection}{\large\sffamily\bfseries\color{ForgeBlue}}{\thesubsection}{0.65em}{}
\titleformat{\subsubsection}{\normalsize\sffamily\bfseries}{\thesubsubsection}{0.65em}{}
\setlength{\parindent}{0pt}\setlength{\parskip}{5.5pt}
\linespread{1.045}
\setcounter{tocdepth}{2}
\setlist{nosep,leftmargin=1.6em,topsep=5pt}
\newcommand{\grind}{\texttt{grind}}
\newcommand{\forge}{\texttt{forge}}
\newcommand{\R}{\mathbb{R}}\newcommand{\Q}{\mathbb{Q}}\newcommand{\N}{\mathbb{N}}
\newcommand{\NN}{\mathbb{N}}
\newcommand{\code}[1]{\texttt{\small #1}}
\newtheorem{proposition}{Proposition}[section]
\newtheorem{lemma}[proposition]{Lemma}
\newtheorem{definition}[proposition]{Definition}
\newtcolorbox{statusbox}{breakable,colback=ForgePale,colframe=ForgeBlue,boxrule=0.4pt,
 left=8pt,right=8pt,top=6pt,bottom=6pt,arc=1pt}
\DefineVerbatimEnvironment{Code}{Verbatim}{fontsize=\small,breaklines=true,breakanywhere=true,
 frame=single,rulecolor=\color{ForgeGray},framesep=6pt,tabsize=2}
\newcolumntype{Y}{>{\raggedright\arraybackslash}X}
\begin{document}
\begin{titlepage}
\vspace*{14mm}
{\sffamily\bfseries\color{ForgeBlue}\fontsize{42}{46}\selectfont FORGE\par}
\vspace{5mm}
{\sffamily\LARGE A Certificate-First Design\\[2mm]Beyond Lean's \texttt{grind}\par}
\vspace{8mm}
{\large Demand-directed proof search, nonlinear certificates,\\
induction, witness synthesis, and checked solver cooperation\par}
\vspace{10mm}
\textcolor{ForgeBlue}{\rule{\textwidth}{0.7pt}}
\vspace{3mm}
{\sffamily\large Research design and executable prototype report\par}
{14 September 2026 \quad / \quad Reference target: Lean 4.34.0\par}
\vspace{7mm}
\begin{abstract}
\noindent
A tactic substantially broader than \grind\ should not replace its congruence
closure with another e-graph, nor merely run a longer list of existing tactics.
It should retain \grind\ as a strong local engine while discovering the missing
intermediate facts, structural arguments, witnesses, and nonlinear certificates
that make local reasoning succeed. This report specifies \emph{Forge}, a typed,
scope-aware, certificate-first architecture for that purpose. Four research
prototypes are implemented and tested: demand-directed Horn inference, sparse
polynomial-cone search, Bernstein positivity certificates, and polynomial
recurrence invariants. They demonstrate concrete mechanisms, including their
failure modes, but do not establish superiority over the actual Lean tactic.
The proposed Lean integration has not been implemented or benchmarked.
\end{abstract}
\vspace{5mm}
\begin{statusbox}
\textbf{Evidence boundary.} The recorded run has 73 passing Python test cases,
109 accepted nonlinear certificates, 43 successful Bernstein certificates out
of 45 cases, and 49 recurrence invariants. Generated Lean replay source is
included, but was \textbf{not compiled or kernel-checked}: no Lean/Lake runtime
was available. This package is a detailed design plus research prototypes,
not a released implementation of \forge.
\end{statusbox}
\vfill
{\small Prepared for the requested investigation of stronger Lean automation.
Source locations, exact dependency versions, raw measurements, tests, and
reproduction scripts accompany the article.}
\end{titlepage}
\setcounter{page}{2}
{\setlength{\parskip}{0pt}\small\tableofcontents}
\clearpage
% BEGIN 01-foundations.tex
\section{The design decision}

\textbf{Build a proof coordinator around \grind, not a replacement for its
successful core.} The proposed tactic, \forge, uses a protected stock-search
lane and a richer search lane. The latter combines backward demands, checked
forward consequences, bounded structural proof search, and specialized
certificate producers. Crucially, those producers can return \emph{useful
lemmas}, not only a verdict on the entire goal. An inequality certificate can
supply an induction step; an induction lemma can expose arithmetic facts;
a synthesized witness can turn a quantified problem into a local \grind\ call.

``More powerful'' has three distinct meanings. \emph{Coverage} concerns which
proof patterns the tactic can discover. \emph{Efficiency} concerns resources
used on common goals. \emph{Reliability} concerns successful replay, predictable
failure, and the trusted assumptions of the resulting theorem. Forge targets
all three, but this report provides experimental evidence only for isolated
Python mechanisms. It does not claim measured Lean-level coverage gains or a
universal speedup.

\subsection{What should become easier}

The intended gains are concentrated where an otherwise short proof needs one
missing idea: strengthening an induction hypothesis, manufacturing a useful
term absent from the ground context, choosing a product of known nonnegative
expressions, finding an existential witness, or extracting a finite Boolean
subproblem. A large library can contain all the necessary lemmas and still
leave these decisions unresolved.

For example, from $x\geq 1$, $y\geq 2$, and $z\geq 3$, a useful intermediate
fact is
\[
 (x-1)(y-2)(z-3)\geq0.
\]
This is not difficult mathematics. The automation problem is selecting that
product, recording its dependencies, and making it available when a polynomial
normalizer or another solver needs it. Likewise, a proof about an accumulator
function may need a universally quantified accumulator invariant before any
amount of ground simplification helps.

The proposed architecture makes those choices explicit, bounded, inspectable,
and replayable. It does not promise automatic discovery of arbitrary creative
mathematics, complete higher-order proof search under finite resources, or
complete nonlinear arithmetic from a fixed-degree certificate basis.

\subsection{Implemented, specified, and still to be built}

\begin{center}
\begin{tabularx}{\textwidth}{@{}p{0.28\textwidth}YY@{}}
\toprule
Component & Delivered evidence & Production work still needed\\
\midrule
Demand-directed inference & Python tabling, dependency closure, proof DAG checker, differential tests & Typed Lean terms, e-class matching, library index, scoped metavariables\\
Nonlinear cone & Exact sparse identities; LP proposals; rational reconstruction & Lean reification, checker theorem or proof replay integration, scheduling\\
Bernstein positivity & Exact box certificates and complete-cover checking & Lean basis lemmas, dependent proof replay, richer split policies\\
Recurrence invariants & Interpolation proposals; exact base and transition identities & Extract Lean recursive definitions and apply appropriate recursors\\
Induction and witnesses & Detailed algorithms and worked derivations & Implement search, generalization, typed candidate synthesis\\
SAT/SMT/superposition & Concrete adapters and trust contracts specified & Version-compatible integrations, proof-rule coverage, testing\\
Whole \forge\ tactic & Architecture, interface, acceptance protocol & Not implemented; no Lean head-to-head results\\
\bottomrule
\end{tabularx}
\end{center}

\section{Begin with the actual baseline}

\subsection{Do not underestimate current \grind}

The current reference describes a cooperating collection of congruence closure,
constraint propagation, E-matching, guided case analysis, and theory solvers.
Its examples include nontrivial polynomial identities and finite-field
reasoning, not merely linear inequalities. It builds ordinary Lean proof
terms. The documentation also identifies large combinatorial Boolean search
as a poor fit and recommends a SAT-backed route for such problems
\cite{lean-grind}.

Consequently, ``add nonlinear arithmetic'' is too imprecise as a contribution.
Forge's algebraic extension is \emph{search for nonnegativity certificates and
useful inequality cuts}, rather than rediscovering polynomial equality
normalization. Similarly, ``add an e-graph,'' ``share facts,'' and ``support
plugins'' describe capabilities already present, not substantive extensions.

The inspected implementation and API also contain theorem guards, delayed
instances, support for extensionality and function congruence,
model-based theory-combination hooks, and a solver-extension interface.
The \code{Action} control layer uses continuations to accommodate branching,
nonchronological backtracking, and replay-script extraction
\cite{grind-types,grind-action}. Forge should exploit those mechanisms instead
of replacing them with an incompatible second notion of proof state.

\subsection{Related systems to reuse rather than relabel}

Aesop supplies a useful model for normalization rules and safe/unsafe
AND--OR search \cite{aesop}. Mathlib's \code{linarith} already separates an
untrusted search oracle from certificate reconstruction; \code{nlinarith}
adds nonlinear preprocessing \cite{linarith}. Thus, an LP oracle followed by an
exact identity check is not itself novel. The extra proposal is a goal-sensitive
product-and-square dictionary, cooperation with structural search, and the
protocol for feeding certified consequences back into a persistent context.

Duper is a concrete candidate for bounded superposition-based reasoning,
including higher-order capabilities \cite{duper}. Lean-SMT demonstrates a
route through preprocessing and proof-producing external SMT search, with
Lean-side reconstruction \cite{lean-smt}. These systems are optional lanes,
not reasons to erase type information or trust an external solver's answer.
A dependency must pass a compatibility and proof-reconstruction test before it
is enabled; the mere existence of a repository does not establish compatibility
with the target toolchain.

Recent work on interventions inside \grind\ is particularly relevant. It
reports that always-on changes can rescue some problems while breaking other
previously successful searches, and advocates interventions after baseline
failure \cite{interventions}. Forge adopts that deployment discipline. It does
not claim that a fallback cascade or bounded lookahead is a new invention.

\subsection{Version and provenance}

The inspected Lean release is \textbf{4.34.0}. The supplied replay project pins
mathlib to commit
\begin{center}
\code{1cf325a0cf67aca2b04d76b5380ff6a9e410aefa}.
\end{center}
The upstream toolchain at that revision specifies Lean 4.34.0. These facts were
read on 14 September 2026 in the America/Los\_Angeles time zone
\cite{lean-release,mathlib-revision}. The corresponding commit timestamp can
fall on 15 September in UTC. Pinning a compatible-looking pair is not a
successful build: the generated project was not compiled here.

\section{Architecture: proof obligations rather than opaque tactic calls}

\subsection{The persistent state}

At a logical level a search node is a sequent $\Gamma\vdash G$, together with a
resource account and a persistent snapshot. Its auxiliary data are:
\begin{itemize}
\item a typed term index and the local \grind\ state, with equality explanations;
\item certified facts, pending guarded facts, and backward demands;
\item an AND--OR graph of proof obligations and alternatives, with dependency
and scope information;
\item a certificate store, learned failure summaries, and a replay-cost estimate.
\end{itemize}

The distinction between \emph{fact} and \emph{demand} is fundamental. A fact
contains an expression $p$ and a proof of $p$ in a specified context. A demand
merely asks for such a proof. A proposed lemma is not inserted into the fact
store until it has been proved. Otherwise, two unsuccessful searches could
``prove'' each other through circular assumptions.

Represent a guarded proposal as
\[
 (\Gamma;\ g_1,\ldots,g_k;\ p;\ c),
 \qquad c:\Gamma\vdash g_1\to\cdots\to g_k\to p.
\]
The watchers for the $g_i$ may awaken when equalities or arithmetic bounds
change. Once each guard has a proof in a permitted scope, ordinary application
of $c$ produces a fact. Before that point, $p$ is only a conditional proposal.
This is also how a division normalization can wait for a nonzero denominator
instead of silently assuming it.

\subsection{Scope, identity, and transport}

Every stored theorem application records its local-context identifier,
free-variable dependencies, instantiated universes, typeclass instances,
metavariable assignments, and proof parents. A fact valid under a branch
assumption cannot leak into its sibling. To move a fact outside its original
scope, first abstract the branch assumptions and then prove the resulting
implication in the receiving context.

For example, a branch may prove $x^{-1}x=1$ using $x\ne0$. After backtracking,
the reusable result is $x\ne0\to x^{-1}x=1$, not the unguarded equality.
Caching only the printed expression loses this distinction. Environment
revision, transparency policy, and dependency closure must be part of a cache
key or explicitly revalidated on lookup.

Lean expressions also cannot be treated as an untyped term algebra.
Homogeneous equality requires compatible types; dependent applications may
need heterogeneous equality and transport. Existing \grind\ services should
supply equality explanations. Forge should never identify terms just because
their pretty-printed strings agree, or because an external first-order encoding
mapped them to the same symbol.

\subsection{The cooperative action protocol}

A plugin returns one of four logical outcomes: no applicable action, checked
progress, a closed proof, or an unresolved result with a reason. A progress
result carries a certified lemma or a well-formed refinement of the goal into
subgoals. Failure, numerical trouble, unsupported syntax, and timeout do not
mean falsehood.

\begin{Code}
cooperate(node, budget):
  saturate cheap certified consequences
  if the goal has a checked proof: return that proof
  enqueue demands extracted from the goal and pending guards
  repeat while resources remain:
    choose an action with a bounded local budget
    run it against an isolated snapshot
    if it returns a certificate:
      reconstruct and check its exact proposition
      merge only scoped, certified consequences
    if it returns subgoals:
      register an AND-node and its proof-combining continuation
    if it closes an OR-alternative:
      retain its proof and cancel unnecessary siblings
    age postponed actions; periodically service the fair lane
  return unknown with residual obligations and resource statistics
\end{Code}

This is stronger than repeatedly restarting \code{first | grind | nlinarith |
aesop}: the plugins exchange selected lemmas, guard proofs, and demand
information. Nevertheless, the first usable integration can begin with
independent leaf calls while the shared-state adapter is tested. That
engineering milestone should not be advertised as the complete architecture.

\subsection{A precise fallback guarantee}

Let $B(G,H)$ denote the result of stock \grind\ on $G$ under a fixed budget
$H$ and a fixed environment. Run that exact computation first. If it succeeds,
return its proof. Only if it fails, restore the initial snapshot and spend an
additional budget $K$ on Forge's extensions.

\begin{proposition}[Additive-budget preservation]
Assuming identical baseline inputs, resource accounting, and execution
conditions, the cascade succeeds on every goal for which $B(G,H)$ succeeds.
\end{proposition}
\begin{proof}
The successful baseline execution is performed unchanged, and its result is
returned before any extension is run.
\end{proof}

This elementary guarantee is about \emph{additional resources}. It is not a
proof that allocating part of the same total budget $H$ to new machinery
preserves all baseline successes. Nor does a simultaneously running baseline
thread receive an identical wall-clock or memory budget if the other threads
contend with it. Report additive-budget and fixed-total-budget experiments
separately. Fairness prevents starvation in an unbounded search; it does not
make a bad finite-budget allocation free.

\section{Demand-directed instantiation and intermediate lemmas}

\subsection{Backward demands with forward proof completion}

Suppose a theorem has an instantiated conclusion matching the wanted fact:
\[
 A_1(\vec t)\to\cdots\to A_k(\vec t)\to C(\vec t).
\]
Rather than forward-generating every instance of every theorem, create demands
for the $A_i(\vec t)$ when $C(\vec t)$ is useful. A conclusion index retrieves
candidate declarations by typed head symbols, argument structure, and relevant
relation classes. Matching is modulo certified equality when useful, but
elaboration decides whether the proposed instantiation is legal.

A candidate whose conclusion does not determine every parameter needs a second
phase: enumerate a bounded typed term grammar, solve a small witness problem,
or leave it to the fair forward lane. It must not arbitrarily instantiate a
universally quantified theorem variable with an ill-scoped metavariable.

The Python prototype deliberately avoids that difficulty. Its rules are
single-sort Horn clauses whose body variables all occur in the head. Matching
a \emph{ground} demanded head therefore determines the substitution for every
premise. There is no higher-order unification, dependent typing, equality
reasoning, or library retrieval in this prototype.

\subsection{Tabling prevents cyclic pseudo-proofs}

Build a finite, budgeted backward closure of ground demands and candidate rule
instances. Then propagate actual proofs forward over that dependency graph.
Each rule instance has a counter of unresolved premises. Input facts initialize
the agenda; proving a premise decreases the counters of its dependents. A rule
fires only when every premise is proved. Store the proof node at that moment.

This is not a depth-first convention that treats ``currently being visited''
as success. Rules $P\to Q$ and $Q\to P$, without a base fact, prove neither.
The certificate is a topologically ordered DAG: every premise index is strictly
smaller than its conclusion's index. Its checker reconstructs the exact
substitution and conclusion for each theorem application.

\begin{proposition}[Relative completeness of the finite Horn core]
For a fixed finite set of ground candidate rule instances, if the dependency
agenda is exhausted without a resource cutoff, every fact in their least
forward closure is eventually proved. In particular, a backward-relevant
closure containing a ground derivation of the target will discover a proof.
\end{proposition}
\begin{proof}
Induct on the height of a finite derivation. Height-zero input facts enter the
agenda initially. Once all smaller-height premises of an instance have entered,
its unresolved counter reaches zero, so its conclusion is enqueued. Tabling
ensures repeated discoveries do not obstruct this propagation.
\end{proof}

The statement concerns the fixed finite closure, not unrestricted first-order
logic. A depth cap can exclude a necessary derivation. With function symbols,
the full relevant closure may be infinite. Such outcomes are reported as
\code{unknown}, not as refutations.

\subsection{A family that isolates the advantage}

Take a base fact $P(a)$ and the two rules
\[
 P(x)\to P(g(x)),\qquad P(x)\to P(f(x)).
\]
Ask for $P(f^d(a))$. In the supplied breadth-first forward implementation, the
$g$-rule precedes the $f$-rule. The goal is therefore the last word generated at
depth $d$, after
\[
 1+2+\cdots+2^d=2^{d+1}-1
\]
facts. Backward demand generation produces exactly the $d+1$ facts on the
$f$-chain. Both returned proofs contain $d+1$ nodes.

At depth 14 the measured counts are 32,767 versus 15. This is an intentionally
adversarial \emph{Python Horn experiment}, not a reconstruction of \grind.
Current \grind\ already sees a negated goal and performs sophisticated
instantiation; this example cannot establish that it would suffer the same
expansion. What the experiment establishes is the implemented dependency
mechanism and its exact work reduction on this family.

\subsection{Library retrieval and induction of useful demands}

A production index should first use deterministic structure: conclusion head,
constants, type constructor, theorem arity, and normalized argument positions.
An inexpensive text ranker can supplement these features. Neural ranking is
optional; it is never proof evidence. After retrieving a declaration, instantiate
it through Lean's elaborator, generate its side goals, and attach its actual
proof term. Never trust a textual theorem statement supplied by a model.

Useful demands need not be the final goal. Sources include a missing denominator
condition, the guard of a conditional rewrite, the precondition of an induction
step, and a sign condition requested by the nonlinear solver. Give these
consumers a small number of ranked requests rather than internalizing every
retrieved theorem globally. Cache failed instantiations by the exact context
and budget, but invalidate failures when relevant facts change.

Demand search, like any relevance filter, can miss a surprising intermediate
lemma. Reserve a fair lane for ordinary instantiation and gradually widen the
retrieval and term-generation bounds. The practical objective is to find a
short proof cheaply, not to turn a relevance heuristic into an unjustified
completeness claim.

% END 01-foundations.tex

% BEGIN 02-algebra.tex
\section{PolyCone: search for nonlinear inequality certificates}

\subsection{Certificate language and its soundness}

Let $p,g_1,\ldots,g_m,f_1,\ldots,f_r\in\Q[x_1,\ldots,x_n]$.
Assume the current context proves $g_i(\vec x)\geq0$ and
$f_j(\vec x)=0$. A certificate for $p(\vec x)\geq0$ is the exact polynomial
identity
\begin{equation}\label{eq:cone}
 p=\sum_{k=1}^{s}\lambda_k q_k^2\prod_{i=1}^{m}g_i^{e_{ki}}
   +\sum_{j=1}^{r}h_j f_j,
 \qquad \lambda_k\in\Q_{\geq0},\quad e_{ki}\in\N,
\end{equation}
where $q_k,h_j$ are rational polynomials. Repetition of a factor is permitted;
the empty product is 1. The equality multipliers $h_j$ may have either sign.

\begin{proposition}[Cone-certificate soundness]\label{prop:cone}
If identity \eqref{eq:cone} holds and all of its sign and equality hypotheses
hold at a real valuation, then $p$ is nonnegative at that valuation.
\end{proposition}
\begin{proof}
Every square is nonnegative. Products of the nonnegative $g_i$, the square,
and $\lambda_k$ are nonnegative, so their sum is nonnegative. Each $h_jf_j$
vanishes at the valuation. Evaluating the polynomial identity gives the claim.
\end{proof}

This proof is small compared with the search for the identity. The implemented
checker independently reconstructs both sums using rational sparse-polynomial
arithmetic, checks factor indices and coefficient signs, and compares the
result with $p$ exactly. It does not consult the LP solver or a numerical
residual threshold. Its input hypotheses remain hypotheses: a Lean adapter
must supply their actual proofs in the current scope.

\subsection{Finite dictionary construction}

For degree cap $D$, generate a bounded collection of columns
\[
 c_k=q_k^2\prod_i g_i^{e_{ki}},\qquad \deg c_k\leq D.
\]
Begin with monomial $q_k$ of degree at most $D/2$. Add binomial candidates
suggested by the target: unit sums and differences, exact rational square-root
ratios of positive even-degree terms, and completing-square expressions
inferred from cross coefficients. If
\[
 a u^2+cuv+bv^2
\]
appears, with $a>0$, try $q=u+(c/(2a))v$; when $b>0$, also try the opposite
orientation. This is a proposal mechanism, not a rule asserting that the
remaining polynomial is nonnegative.

The completing-square proposal matters in practice. The initial heuristic
recognized $(x-2/3)^2$, but a positive extra constant obscured its coefficient
ratio in $(x-2/3)^2+1/97$. Adding the cross-coefficient construction recovered
the intended square. A regression test now covers this case. The unsuccessful
initial search returned \code{unknown}; it did not accept an approximate proof.

Form products of known nonnegative constraints subject to the same degree cap.
Include columns $m f_j$ for monomials $m$ of degree at most $D-\deg f_j$; their
multipliers are unrestricted rational variables. Deduplicate identical columns
and retain their construction recipes. Avoid expanding the global Cartesian
product when most constraints concern unrelated variables: partition by
variable support and request cross-component products only when the target
connects those components.

The shipped implementation uses a simple bounded enumeration and a limit of
4,000 generators. It is a research implementation, not an optimized sparse
SOS package. A production implementation needs a separate bound on enumeration
work, monomial count, coefficient bit length, and external-solver time; a cap
on successfully admitted columns alone is insufficient as a denial-of-service
or worst-case complexity control.

\subsection{Search becomes linear algebra with an untrusted oracle}

Expand the columns and target in a common monomial basis. The certificate
search is
\begin{equation}\label{eq:lp}
 A\lambda+Bh=b,\qquad \lambda\geq0,
\end{equation}
where the $h$ coordinates are free. The Python prototype calls SciPy's
\code{linprog(method="highs")} to propose a solution \cite{scipy-lp}.
It minimizes the sum of cone coefficients as a simple heuristic. A production
objective should also penalize large polynomial terms, expensive guard proofs,
and predicted replay cost; minimizing coefficient sum is not proof-size
minimization.

Floating-point coefficients are converted to rational candidates with bounded
denominators and checked exactly. If that fails, solve the proposed active
support with exact SymPy linear algebra, set any remaining free parameters to
zero, and check the result again. A negative reconstructed coefficient, an
inconsistent support, or a failed identity causes rejection. No tolerance is
used at the acceptance boundary.

\begin{Code}
polycone(target, inequalities, equalities, D):
  columns := degree-bounded products and proposed squares
  add ideal columns with unrestricted coefficients
  numerical_candidate := solve_linear_program(columns, target)
  if no candidate: return unknown
  rational_candidate := reconstruct(numerical_candidate)
  if exact_check(rational_candidate): return certificate
  support_candidate := exact_solve(numerical_support)
  if exact_check(support_candidate): return certificate
  return unknown
\end{Code}

The choice of zero values for free reconstruction parameters is incomplete:
a different rational point on the same support could have nonnegative
coefficients. Similarly, numerical infeasibility does not establish
nonnegativity failure and may not even establish infeasibility of the exact
finite problem. These are performance limitations only because every
successful result passes an independent exact check.

\subsection{Worked certificates}

Several tested cases explain the range of the certificate language.
For the two-dimensional Cauchy--Schwarz inequality,
\[
 (a^2+b^2)(c^2+d^2)-(ac+bd)^2=(ad-bc)^2.
\]
For variance in three variables,
\[
 x^2+y^2+z^2-xy-yz-zx
 =\tfrac12(x-y)^2+\tfrac12(y-z)^2+\tfrac12(z-x)^2.
\]
Both follow from squares without order assumptions on the variables.
For $x,y\in[0,1]$,
\[
 xy(1-x)(1-y)\geq0
\]
uses four known nonnegative factors. The identity has degree 4 and is rejected
by the prototype's degree-2 search before LP construction.

Equality hypotheses contribute differently. Under $y=x$,
\[
 y^2-2x+1=(x-1)^2+(x+y)(y-x).
\]
The second term need not be nonnegative; it vanishes because of the equality.
Treating all certificate multipliers as nonnegative would unnecessarily lose
this useful algebraic flexibility.

\subsection{Theory cooperation through generated inequalities}

The plugin should return useful consequences before it can prove the final
goal. Suppose $l_x\leq x\leq u_x$, $l_y\leq y\leq u_y$, and a shared atom
$z$ is certified equal to $xy$. The four products of interval distances yield
\begin{align*}
 z&\geq l_x y+l_y x-l_xl_y, &
 z&\geq u_x y+u_y x-u_xu_y,\\
 z&\leq u_x y+l_y x-u_xl_y, &
 z&\leq l_x y+u_y x-l_xu_y.
\end{align*}
For instance, the first follows from $(x-l_x)(y-l_y)\geq0$.
These inequalities are linear in the shared atoms $x,y,z$ and can strengthen
an existing linear arithmetic solver. Their derivations remain polynomial
certificates with explicit bound dependencies. The proposal is to generate
such cuts on demand, not to replace the linear engine.

For natural numbers and integers, use their actual semantics. Natural
subtraction is truncated; machine arithmetic may wrap; division and remainder
have domain-specific rules. Reify into an ordered field only after proving the
necessary cast and operation correspondence. A real-arithmetic relaxation can
supply a sufficient certificate for an integer theorem after a checked
embedding, but it cannot justify an integer witness or an integer-specific cut
without an additional argument.

Strict inequalities require separate bookkeeping. A positive rational margin
in a certificate is enough to prove strict positivity; nonnegativity alone is
not. Rational-expression normalization requires denominator conditions. These
conditions are obligations on the fact graph, not informal side notes in an
external certificate.

\subsection{Limits and a concrete SOS escalation}

The number of monomials of total degree at most $D$ in $n$ variables is
$\binom{n+D}{D}$. Even before proof reconstruction, indiscriminate degree and
variable growth can make the coefficient problem too large. Apply cheap
symbolic factorization and variable separation first, then expand a bounded
basis in increments, recording which expansion helped.

A finite binomial-square dictionary is strictly more restrictive than arbitrary
sum-of-squares search. An optional escalation can introduce a Gram matrix
$Q\succeq0$ with $p=v^{\mathsf T}Qv$ or its constrained analogue, solve for $Q$
using a semidefinite optimizer, and reconstruct an \emph{exact} rational
factorization
\[
 Q=L\,\operatorname{diag}(d_1,\ldots,d_t)L^{\mathsf T},\qquad d_i\geq0.
\]
The resulting weighted squares fit the same certificate checker. A numerical
claim of positive semidefiniteness is not accepted; an exact factorization and
identity are required. Singular matrices and rational reconstruction can make
this escalation fail. No SDP solver is implemented in this package.

The tested Motzkin polynomial
\[
 M(x,y)=x^4y^2+x^2y^4+1-3x^2y^2
\]
is nonnegative by arithmetic--geometric mean applied to the three nonnegative
terms $x^4y^2$, $x^2y^4$, and 1. The finite unconstrained square dictionary does
not certify it. Its status is therefore \code{unknown}, despite truth. No
bounded cone in this report is advertised as a complete real-algebraic decision
procedure. Rational multipliers, richer positivity theorems, or other proof
methods require explicit new certificate forms and soundness lemmas rather
than a more generous floating-point tolerance.

\section{Bernstein certificates on rational boxes}

\subsection{A complementary proof language}

A local box constraint can make a positivity proof easier than a global
square decomposition. Map $x\in[a,b]$, $a<b$, to $t=(x-a)/(b-a)\in[0,1]$.
For degree $d$ write the polynomial exactly as
\[
 p(a+(b-a)t)=\sum_{k=0}^{d} b_k B_{k,d}(t),\qquad
 B_{k,d}(t)=\binom{d}{k}t^k(1-t)^{d-k}.
\]
If every $b_k\geq0$, each summand is nonnegative on the interval. Products of
these bases give a certificate for a multivariate box.

For power coefficients $a_j$ on $[0,1]$, the conversion formula is
\[
 b_k=\sum_{j=0}^{k}a_j\frac{\binom{k}{j}}{\binom{d}{j}}.
\]
To verify it, use
$t^j=\sum_{k=j}^{d}\binom{k}{j}\binom{d}{j}^{-1}B_{k,d}(t)$,
which follows by expanding the binomial theorem for the remaining degree
$d-j$. The multivariate conversion is its tensor product.

The prototype discovery routine computes the coefficients after an affine
change of variables. Its checker takes a deliberately different route:
it reconstructs every basis term directly in the original coordinates,
\[
 \binom{d}{k}\frac{(x-a)^k(b-x)^{d-k}}{(b-a)^d},
\]
and checks the resulting polynomial identity exactly. Both routines share the
sparse rational arithmetic layer, so their implementation independence is not
complete. Differential tests against SymPy address that shared layer but do
not constitute formal verification.

\subsection{Subdivision is a proof tree, not sampling}

If some coefficients are negative, bisect a box coordinate and try both
children. A successful certificate is a complete binary tree of boxes, with
nonnegative-coefficient leaves. The checker verifies the root box, each cut's
location, the exact two child boxes, and every leaf identity. A gap in the
cover invalidates the entire certificate even if every retained leaf looks
correct. Testing finitely many points is never an acceptance rule.

For
\[
 p(x)=x^2-x+\tfrac13\quad\text{on }[0,1],
\]
the degree-2 coefficients are $(1/3,-1/6,1/3)$, so the initial box does not
certify positivity. On $[0,1/2]$ they become $(1/3,1/12,1/12)$, and on
$[1/2,1]$ they become $(1/12,1/12,1/3)$. Two leaves prove the entire interval.
This example is included in the tests and measurements.

\subsection{Why true statements can remain unresolved}

The polynomial $(x-1/3)^2$ is nonnegative, yet midpoint subdivision never makes
$1/3$ an endpoint of an interval with dyadic endpoints. On the leaf containing
that point in its interior, all Bernstein basis functions are strictly
positive there. If every coefficient were nonnegative and their sum were zero
at that interior point, every coefficient would have to vanish. The polynomial
would then vanish identically on the leaf, which it does not. Thus this
particular subdivision strategy cannot finish, not merely because the tested
depth cap is too small.

This is a useful division of labor: the cone routine can recognize the square,
while a naive box routine cannot. Conversely, box constraints can certify
polynomials that a small global dictionary misses. Route a problem according
to observed structure, preserve both alternatives, and report the actual
certificate family that succeeded.

A production splitter can use rational root isolation or symbolic factorization
to expose zero sets before subdivision. That is a concrete improvement, but
it requires certified root information and complete domain decomposition.
Adaptive numerical guesses about root locations remain proposals until those
obligations are discharged.

% END 02-algebra.tex

% BEGIN 03-structure.tex
\section{Induction as a search action}

\subsection{Generate a justified induction problem}

More ground reasoning is not a substitute for a structural argument. Forge
should inspect recursive definitions occurring in the goal, their equation
theorems, and the arguments on which recursive calls decrease. This yields a
small, ranked set of induction candidates. Structural induction on a list or
natural number is the first case; a function-specific induction theorem or a
well-founded recursor is used when the recursion requires it.

The algorithm must preserve the distinction between discovering an induction
hypothesis and assuming an arbitrary stronger assertion. A generalization
creates a \emph{new theorem to prove}, with its own base and step obligations.
Only the checked induction principle closes that theorem. An open induction
branch cannot contribute its conjectured invariant as a global fact.

A concrete procedure is:
\begin{enumerate}
\item Read the target and selected hypotheses as a dependency graph. Identify
recursive calls and the local variables in their changing argument positions.
\item Rank induction variables using occurrence in decreasing arguments,
recursor size, and expected branch count. Start with one variable, then bounded
pairs or a function-specific induction principle.
\item Form a generalization set from variables changed by the recursive calls.
Take its dependent-type closure, so reverting a variable does not strand local
declarations whose types depend on it.
\item Revert the selected declarations, apply the actual recursor or induction
theorem, introduce branch parameters and induction hypotheses, and unfold only
relevant equation lemmas.
\item Solve the resulting AND-branches cooperatively. On failure, retain the
residual goal to suggest a wider generalization or a missing auxiliary lemma,
but do not reuse the unproved assertion as evidence.
\end{enumerate}

A ``safe'' transformation is not necessarily cheap. Even semantically
appropriate induction can create an enormous proof search. Put a hard bound
on branch count and total candidate recursors before invoking it.

\subsection{Worked accumulator example}

Define, mathematically,
\begin{align*}
 \operatorname{revAcc}([],a)&=a,\\
 \operatorname{revAcc}(x::s,a)&=\operatorname{revAcc}(s,x::a).
\end{align*}
The natural target is
$\operatorname{revAcc}(s,[])=\operatorname{reverse}(s)$.
A weak induction hypothesis talks only about the accumulator $[]$, whereas the
recursive call uses $x::a$. It does not match.

The changed-argument analysis identifies the accumulator, suggesting
\begin{equation}\label{eq:rev-inv}
 \forall a,\quad \operatorname{revAcc}(s,a)
       =\operatorname{reverse}(s)\mathbin{+\!+}a.
\end{equation}
The empty-list case follows from the equations. In the cons case, the induction
hypothesis instantiated at $x::a$ gives
\[
 \operatorname{revAcc}(x::s,a)
 =\operatorname{reverse}(s)\mathbin{+\!+}(x::a).
\]
The recursive equation for reverse and associativity of append identify this
with $\operatorname{reverse}(x::s)\mathbin{+\!+}a$.
Specializing the now-proved generalized theorem at $a=[]$ gives the original
target. The complete proof idea is the quantified accumulator invariant, not
an unexplained instruction to ``try induction harder.''

The library may already contain a theorem that makes this example trivial.
That does not invalidate the generalization mechanism, but it changes the
benchmark. Tests of structural synthesis must use fresh definitions or control
which equivalent high-level lemmas are available. Never count retrieval of an
already-proved target as evidence that induction synthesis worked.

\subsection{Failure-guided auxiliary lemmas}

After an unsuccessful step, compare the recursive call with the induction
hypothesis. Mismatched accumulator arguments suggest generalization;
mismatched compositions suggest a fusion lemma; polynomial differences suggest
an invariant template. A candidate lemma is ranked by term size, number of new
quantifiers, and whether its recursive calls have a clear decreasing argument.
Try at most a bounded number at each size.

For a function $f$ with an accumulator update $a\mapsto T(a,x)$, templates can
include
\[
 f(s,a)=H(s,a),\qquad
 \mu(f(s,a))=\mu(a)+J(s),\qquad
 R(f(s,a),s,a).
\]
The template vocabulary is built from the goal, equation theorems, and a small
approved operator set. The resulting lemma is an OR-alternative with a separate
AND proof tree. If its proof closes, inject its specialization into the parent.
Otherwise discard it or widen the budget. This avoids mutually recursive
``helper lemmas'' becoming a circular proof.

\section{Polynomial invariant synthesis with exact validation}

\subsection{From recurrence data to a theorem candidate}

For a state transition $T$ and starting state $s_0$, an equality invariant
$I(s)=0$ is justified by
\[
 I(s_0)=0,\qquad I(T(s))=I(s).
\]
The second identity is a convenient sufficient condition. More generally it
would suffice to prove $I(s)=0\Rightarrow I(T(s))=0$, but that is not the
certificate language implemented here.

The prototype considers additive polynomial recurrences
\begin{equation}\label{eq:additive}
 n'=n+1,\qquad s'=s+p(n),\qquad (n_0,s_0)=(0,c),
\end{equation}
where $p\in\Q[n]$. It proposes a polynomial $Q(n)$ and sets
$I(n,s)=s-Q(n)$. The certificate obligations become
\begin{equation}\label{eq:fdiff}
 Q(0)=c,\qquad Q(n+1)-Q(n)=p(n).
\end{equation}
Given a degree bound $d$, sampling the recurrence at $d+1$ points and
interpolating suggests $Q$. \emph{Sampling is only discovery.} The validator
substitutes the polynomial transition into $I$ and compares exact coefficient
maps for the universal identity.

A production synthesizer should also support direct coefficient solving.
Write $Q(n)=\sum_{i=0}^{d}c_i n^i$, expand \eqref{eq:fdiff}, and solve the
resulting rational linear system. This removes numerical concerns and avoids
a dependence on the availability of inexpensive evaluation samples. The
interpolation implementation was chosen here because it makes the separation
between plausible evidence and universal checking particularly explicit.

\subsection{Example: sum of squares}

For $p(n)=(n+1)^2$ and $c=0$, the discovered expression is
\[
 Q(n)=\frac{n(n+1)(2n+1)}6.
\]
Equivalently take the denominator-cleared invariant
\[
 J(n,s)=6s-n(n+1)(2n+1).
\]
Its step difference is
\begin{align*}
 J(n+1,s+(n+1)^2)-J(n,s)
 &=6(n+1)^2\\
 &\quad-(n+1)(n+2)(2n+3)+n(n+1)(2n+1)\\
 &=0.
\end{align*}
Its base value is zero. Ordinary induction on the number of iterations
therefore proves the equality for every reachable state. No extrapolation from
finite data is involved after the two identities have been checked.

The tests include a hostile interpolation example: an added high-degree term
vanishes at all sampled points but violates the recurrence identity elsewhere.
The exact validator rejects the proposed invariant. This is an essential test
for any counterexample-guided or sample-based synthesis component.

\begin{proposition}[Invariant certificate soundness]
Let $s_{k+1}=T(s_k)$. If $I(s_0)=0$ and $I\circ T=I$, then
$I(s_k)=0$ for every $k\in\N$.
\end{proposition}
\begin{proof}
The claim for $k=0$ is the base condition. If $I(s_k)=0$, then
$I(s_{k+1})=I(T(s_k))=I(s_k)=0$. Apply induction.
\end{proof}

\subsection{Generalizing the invariant engine}

For a polynomial state vector $\vec x$, use the template
\[
 I(\vec x)=\sum_{|\alpha|\leq d}c_\alpha\vec x^\alpha.
\]
Because $T$ is fixed, the coefficients of $I(T(\vec x))-I(\vec x)$ are linear
in the unknown $c_\alpha$. A rational nullspace computation discovers conserved
polynomials, with the additional linear constraint $I(s_0)=0$. Reject the
identically zero invariant and rank remaining candidates by relevance to the
postcondition. This is an immediately implementable extension, although the
shipped discovery routine handles only \eqref{eq:additive}.

A more expressive certificate permits a vector of invariant equalities:
\[
 I_i(T(\vec x))=\sum_j h_{ij}(\vec x)I_j(\vec x).
\]
This proves preservation of their common zero set. Bounded-degree multiplier
search reduces the check to exact identities, but jointly discovering the
$I_j$ and multipliers is no longer the same easy linear problem. Separate
candidate generation from multiplier certification, or fix one side at a time;
do not describe the joint bilinear search as linear algebra.

Inequality invariants require base and preservation \emph{inequalities}. The
PolyCone plugin can certify these under transition guards and existing
invariants. Termination remains a different obligation: an invariant does not
prove that a loop or recursive definition terminates. For already accepted
Lean definitions, use their existing recursion principles; for a proposed new
recursive program, termination must be proved separately.

\section{Existential witnesses and counterexample-guided synthesis}

\subsection{Construct witnesses in the correct scope}

For a goal $\exists y:\tau,\ P(\vec x,y)$, the useful output is a term $t$ of
\emph{exactly} type $\tau$ and a proof of $P(\vec x,t)$. Search begins with
terms already in scope, constructor applications, affine expressions suggested
by equations, and terms built from the relevant theorem signatures. Candidate
terms must not depend on variables introduced inside a later branch or a
future universal quantifier.

For
\[
 \forall n\in\N,\quad\exists m\in\N,\quad m>n\ \wedge\ m\equiv0\pmod2,
\]
a bounded affine grammar can propose $m=2(n+1)$. The remaining proof obligations
are natural-number arithmetic and divisibility. Construct the existential
proof explicitly after those obligations close. A real-valued LP solution for
$m$ would not be a natural-number witness; domain checks are part of the
construction, not an optional verification pass.

\subsection{A concrete CEGIS loop}

Counterexample-guided inductive synthesis (CEGIS) alternates candidate proposals
and attempted refutations of their universal correctness. For an affine
integer grammar $t(\vec x)=c_0+\sum_i c_i x_i$, enumerate bounded integer
coefficient vectors or ask an integer solver for coefficients that satisfy a
finite sample set. Once a candidate is found, ask the cooperative prover to
establish $\forall\vec x,\ P(\vec x,t(\vec x))$ under the actual assumptions.

\begin{Code}
synthesize(grammar, samples, goal, budget):
  repeat within term-size and coefficient bounds:
    t := candidate fitting every validated sample
    result := prove_universally(goal specialized at t)
    if result is a checked proof:
      return explicit witness t and proof
    if result contains a validated counterexample:
      add it to samples
    otherwise:
      schedule a different candidate or a larger proof budget
  return unknown
\end{Code}

An external model is not automatically a Lean counterexample. Recheck its
values against the encoded operations and assumptions, especially casts,
truncating subtraction, division, and finite types. A spurious model should
refine the encoding or be ignored; it must not eliminate a valid candidate.
Likewise, failure to prove a candidate does not mean it is false.

This module is specified, not implemented in the package. A reasonable first
implementation restricts the grammar to constructors, affine arithmetic, and
finite tuples, reusing \code{omega}, \code{linarith}, and \grind\ for residual
proofs. The current reference explicitly does not describe \code{omega} as a
complete procedure for every integer-linear problem \cite{lean-tactics}; the
integration should preserve ordinary tactic failure semantics.

\subsection{Models guide search; they do not certify theorems}

The same discipline applies to learned premise selection, algebraic pattern
recognition, and suggestions from a computer algebra system. A proposed
factorization becomes useful only after the exact identity is proved. A
proposed induction invariant becomes useful only after base and preservation
are proved. A proposed witness becomes useful only after its universally
scoped specification is proved. This uniform boundary makes the architecture
compatible with more ambitious search methods without enlarging the logical
trust base merely because a search component is complicated.

% END 03-structure.tex

% BEGIN 04-engineering.tex
\section{Certified solver islands and theory combination}

\subsection{Export a fragment, not an untyped approximation of Lean}

A \emph{solver island} is a selected subproblem whose translation and proof
reconstruction are supported end to end. Boolean and fixed-width bit-vector
islands can use the existing SAT-oriented facilities. A suitable first-order
island can use Lean-SMT and its reconstruction path; a clause-based island can
use Duper \cite{lean-grind,lean-smt,duper}. The island may prove a lemma over
shared interface terms rather than the complete original goal.

Each export has a translation ledger: Lean expressions, their solver symbols,
their types and operations, and proof obligations connecting the encoding to
the original proposition. Universe instantiation, monomorphization, typeclass
resolution, and auxiliary assertions must have justified transformations.
The Lean-SMT work illustrates why that preprocessing is a substantial part of
an integration rather than mere serialization \cite{lean-smt}.

For an island under assumptions $A_1,\ldots,A_k$, reconstruct a proof of
$A_1\to\cdots\to A_k\to L$. Apply it to the appropriate scoped proofs to
insert $L$ into the main graph. An external \code{unsat} answer without a
supported proof is \code{unknown}. An external \code{sat} answer is a model
proposal, not necessarily a counterexample to the Lean goal. Unsupported proof
rules and unsupported theories must be detected and reported, not silently
trusted.

\subsection{Do not import unsound combination assumptions}

A generic textbook theory-combination procedure is not automatically valid
for every Lean type. Finite domains, dependent types, casts, and overlapping
arithmetic operators complicate disjointness and cardinality assumptions.
Forge avoids relying on a blanket combination theorem: plugins exchange
\emph{proved interface facts}, and every final step is checked in Lean.
This is sound even when incomplete.

When a plugin wants $u=v$ to reconcile a model, it may request a case split or
ask another plugin to prove it. It cannot simply merge those terms. Similarly,
an arithmetic atom that abbreviates $xy$ must carry a definition or equality
proof; treating it as an unrelated real variable gives only a relaxation, not
a bidirectional equivalence of problems.

\subsection{Two trust profiles}

The default target is \textbf{kernel-checked replay}: ordinary proof terms or a
proved certificate-checker theorem whose acceptance computation is justified
without trusting an external search program. The intended trusted assumptions
are Lean's ordinary accepted foundations and explicitly audited library
axioms. Replay should not introduce \code{sorryAx}, a newly assumed theorem,
or a hidden oracle axiom.

A separate, explicitly named \textbf{native-checking profile} may run a
certificate checker using native compilation for performance. This has a
different trust story. The reference warns that \code{native\_decide} adds the
Lean compiler to the trusted part, whereas \code{decide\_cbv} does not require
trust in the code generator \cite{lean-tactics}. Do not summarize both profiles
as ``only the kernel is trusted.'' A SAT certificate and a theorem about its
checker are not, by themselves, a complete description of how the checker was
actually evaluated.

The precise existing implementation of a delegated tactic must be audited
under the pinned version. In strict mode, disable a backend when its available
replay path requires the native profile. Certificate size, kernel-replay time,
and axiom dependencies belong in the result record alongside search time.

\section{Concrete Lean integration plan}

\subsection{Use the current extension interface}

The inspected API has \code{Lean.Meta.Grind.SolverExtension} and
\code{registerSolverExtension}. The latter registers an extension during
initialization and allocates its state. Its methods include internalization,
new equality and disequality notifications, an action, consistency-related
checks, and a model-combination hook. State access, term marking, fact insertion,
and proof explanations are also exposed \cite{grind-types}.

The API skeleton below records actual conceptual hook names; the handler
bodies and Forge state are \textbf{design pseudocode}, not a compiled module.

\begin{Code}
initialize polyExt := registerSolverExtension makePolyState
initialize polyExt.setMethods
  (internalize := collectSupportedArithmetic)
  (newEq       := updateCanonicalAtoms)
  (newDiseq    := wakeNonzeroGuards)
  (action      := boundedCertificateAction)
  (checkInv    := checkInternalBookkeeping)
\end{Code}

Make the adapter a narrow, version-specific module. It reifies supported
expressions with a valuation map, requests a certificate, reconstructs a
proof, and submits a proved fact through the appropriate insertion service.
The check of internal data-structure invariants is not the mathematical
soundness proof of a certificate. Name and test those two responsibilities
separately.

Structural induction is better implemented above the local saturation loop:
it changes the goal shape and creates dependent subgoals. The resulting branch
contexts can each host a \grind\ session and its plugin state. This two-level
architecture keeps the existing theory engine local while the outer graph
tracks recursor applications and generalized lemmas.

\subsection{Module boundaries and required contracts}

\begin{center}
\begin{tabularx}{\textwidth}{@{}p{0.27\textwidth}YY@{}}
\toprule
Proposed module & Responsibility & Acceptance contract\\
\midrule
\code{Forge.Context} & Scoped expression, fact, and obligation graph & No unproved facts; valid dependency scopes\\
\code{Forge.GrindAdapter} & Version-specific extension hooks & Restorable state; exact expression/proof transport\\
\code{Forge.Demand} & Retrieval, instantiation, tabling & Every application elaborates and has acyclic proof parents\\
\code{Forge.Poly.Reify} & Typed polynomial representation & Denotation equality for each reified term\\
\code{Forge.Poly.Cert} & Cone and box certificate checking & Proved checker soundness or explicit proof reconstruction\\
\code{Forge.Induct} & Recursor selection and generalization & Goal refinement through an actual induction theorem\\
\code{Forge.Synth} & Invariants and existential candidates & Base/step or witness specification discharged\\
\code{Forge.Islands} & SAT, SMT, and superposition adapters & Supported translation and proof rules; no raw oracle acceptance\\
\code{Forge.Replay} & Proof minimization and audit & Exact original goal, no remaining metavariables, approved axioms\\
\bottomrule
\end{tabularx}
\end{center}

Reification deserves particular attention. To use a rational polynomial
certificate, store a map $\rho$ from polynomial variables to Lean expressions
and a proof that the original term equals the polynomial's denotation at
$\rho$. For an ordered-field inequality, every coefficient cast and operation
must agree with that denotation. An algebraically correct certificate for the
wrong denotation is not a proof of the original statement.

\subsection{Libraries and algorithms with specific roles}

The initial Lean dependency should be \textbf{core Lean plus mathlib}, with
Aesop used through the version selected by the pinned dependency graph when
available. Reuse \code{ring} or \code{ring\_nf} for polynomial identities,
\code{norm\_num} for rational constants, \code{positivity} for cheap sign facts,
and \code{linarith}/\code{nlinarith} for arithmetic leaves and selected
certificate reconstruction \cite{ring,positivity,linarith}. Do not invoke all
of them blindly on every node.

The research implementation uses Python's \code{fractions.Fraction} and a
small sparse-polynomial representation for checking, SciPy/HiGHS for LP
proposals, and SymPy for exact support solving and interpolation. Exact tested
versions are recorded with the results. NumPy constructs numerical matrices;
its floating-point values never decide certificate acceptance.

Optional integrations are Duper for superposition and Lean-SMT with its
supported external solver for SMT proofs. Keep them out of the minimal build
until compatibility tests pass. A numerical semidefinite solver can later
supply Gram-matrix proposals, but an exact replay path is a prerequisite.
An existing e-graph package is not required: \grind\ already maintains
congruence machinery. In particular, an archived experimental Lean equality
saturation project should not be made a critical dependency simply because
its name sounds aligned with the design.

\subsection{Resource allocation and proof-cost awareness}

After the protected baseline attempt, a practical starting policy allocates
local bounded actions rather than unbounded calls: demand expansion, a small
cone, a shallow induction candidate, a small witness grammar, and an eligible
solver island. Rank actions by observed recent payoff, target relevance, cost
estimate, and age. Reserve a periodic fair slot for postponed applicable
actions. The precise weights are tunable engineering parameters, not
probabilistically calibrated guarantees.

For an implementable non-learned initial scheduler, use an eight-slot round
with three demand slots, two cone slots, one structural slot, one witness slot,
and one eligible external-island slot. Empty slots are donated to the oldest
applicable action. Within a slot, rank target-linked obligations before generic
ones, then prefer lower estimated cost; use a deterministic identifier to break
ties. Every fourth round give the oldest postponed action priority. This is a
proposed starting policy, not a measured optimum or a finite-budget guarantee.

Initial per-action limits can be explicit: expand at most 32 theorem instances
per demand quantum; try cone degrees 2, 4, and 6 with at most 4,000 columns;
try at most two induction candidates per structural quantum and induction depth
at most two; enumerate at most 256 witness proposals per quantum with term size
at most six. Track external jobs separately under wall-clock and output-size
limits. Every quantum also consumes the common parent budget, so nested actions
cannot multiply the allowed work by repeatedly resetting their local limit.

Keep a monotone context epoch alongside each immutable exported problem. An
asynchronous result is reconstructed only after checking that its expression
mapping and hypothesis dependencies still apply. Never reuse an old mutable
metavariable assignment merely because its printed goal matches the current
one. If transport to the new context cannot be justified, discard the result
or rerun the action. These rules are especially important when induction,
backtracking, and external jobs overlap.

Count the cost of \emph{reconstruction}. An apparently fast oracle that emits
an enormous proof can lose to a slower search with a small certificate. Use
support minimization, common subexpression sharing, and proof-DAG trimming.
Keep unsat-core or used-premise information as a proposal for minimization,
then replay the reduced proof to verify that no necessary dependency was
removed. Hash-consing expressions reduces duplication but does not excuse
missing scoping checks.

All external requests require limits on time, memory, serialized input/output
size, integer bit length, and parser depth. Do not evaluate returned text as
code. A well-formed but gigantic certificate is a resource risk even when its
mathematics is harmless. Permit cancellation between checker steps and return
an unresolved result on resource exhaustion.

\subsection{User-facing syntax and explanations}

The following is \textbf{proposed syntax}, not syntax accepted by the delivered
package:
\begin{Code}
by
  forge (mode := conservative)
    (inductionDepth := 2)
    (polyDegree := 4)
    (witnessSize := 6)
    (external := false)
\end{Code}

A companion \code{forge?} should suggest a stable replay script or an explicit
certificate invocation. Users should see which premises were used, which
induction variable was selected, which parameters were generalized, and what
witness was constructed. Failure diagnostics should distinguish a missing
sign guard, a degree cap, unsupported reification, an unproved generalized
lemma, an unavailable backend, and an exhausted search budget.

The report's generated Lean files are a narrower first artifact: explicit
proofs for saved certificates. They contain no search implementation of
\forge. Hiding that distinction behind an umbrella import would misrepresent
the state of the project.

\section{Soundness obligations for the whole system}

\begin{proposition}[Compositional soundness target]
Suppose every fact node contains a Lean proof in its recorded context; every
refinement node is justified by an accepted Lean theorem or recursor; every
certificate adapter proves its exact translated conclusion; and every scope
transfer is checked. Then a closed root of the acyclic proof graph supplies a
proof of the original goal, independently of how search candidates were
selected.
\end{proposition}
\begin{proof}
Proceed in topological order through the proof graph. Base facts are valid by
assumption. At a theorem-application node, apply its proof to the previously
constructed premise proofs. At a refinement node, combine the checked branch
proofs with its theorem or recursor. At a certificate node, use the adapter's
proved conclusion. Scope transfer abstracts and instantiates only validated
dependencies. The final node therefore inhabits the original proposition.
\end{proof}

This is a design-level argument, not a mechanized proof of the present Python
implementation. Lean integration must instantiate each hypothesis of the
argument. An unsound term translator is not repaired by an exact polynomial
checker. A sound checker is not repaired by a replay path that assumes its
result through an oracle axiom. An acyclic certificate DAG is not sufficient
if a proof parent was imported from the wrong branch.

Testing should deliberately attack those boundaries. Mutate a denominator
condition; swap two polynomial variables; replace a strict inequality with a
weak one; change an exponent; reverse an equality transport; remove a branch
from a box cover; alter a typeclass instance; point a proof node to a future
node; and replay under a different local context. The intended result is
rejection, not a best-effort reinterpretation of the malformed evidence.

% END 04-engineering.tex

% BEGIN 05-experiments.tex
\section{Experiments actually performed}

\subsection{Environment and reproducibility}

The execution environment was Linux x86\_64 with Python 3.13.5, NumPy 2.3.5,
SciPy 1.17.0, SymPy 1.14.0, and pytest 9.0.2. The synthetic benchmark uses seed
20260914. Numerical-library thread environment variables were set to one for
the recorded run. Full version strings are in \path{results/benchmarks.json}.
The test run completed with \textbf{73 passed in 1.11 seconds}.

No \code{lean}, \code{lake}, or successful Lean execution route was available.
The supplied Lean comparison runner therefore recorded \code{not\_run}, with
that reason, rather than manufacturing zero failures or baseline timings.
Source inspection, Python execution, and generated Lean source are three
different forms of evidence and remain separately labeled throughout the
package.

The benchmarks were executed once for diagnostic purposes; there are no
repeated-trial confidence intervals or hardware-normalized performance claims.
The exact combinatorial counts and algebraic acceptance results are more
informative here than submillisecond timings.

\subsection{Summary of observed outcomes}

\begin{center}
\begin{tabularx}{\textwidth}{@{}p{0.34\textwidth}rY@{}}
\toprule
Experiment & Outcome & Interpretation\\
\midrule
Pytest suite & 73 passed & Includes parametrized, randomized, negative, and mutation tests\\
Named nonlinear cases & 9 accepted / 12 & Two true cases unresolved; one deliberately false case unresolved\\
Generated product-cone cases & 100 / 100 & Constructed inside the covered cone; not representative problem sampling\\
Separate false-polynomial controls & 0 accepted / 24 & All returned unknown; exact negative witnesses checked separately\\
Bernstein cases & 43 accepted / 45 & One true nondyadic-zero case and one false case unresolved\\
Additive recurrence invariants & 49 / 49 & Polynomial increments; base and universal step identities checked\\
Horn growth family & 5 depths & Both strategies prove all goals; search work differs sharply\\
Actual Lean comparison & Not run & No measured superiority over \grind\\
\bottomrule
\end{tabularx}
\end{center}

The 109 accepted nonlinear certificates are the nine named successes plus the
100 generated cases. They are not 109 previously unsolved Lean theorems. The
49 invariants concern a deliberately restricted polynomial recurrence class,
not arbitrary recursive Lean programs.

\subsection{Demand search: counts before speed claims}

% BEGIN horn-table.tex
\begin{center}
\begin{tabular}{@{}rrrrr@{}}
\toprule
Depth & Forward facts & Demand atoms & Forward ms & Demand ms\\
\midrule
4 & 31 & 5 & 0.416 & 0.128 \\
8 & 511 & 9 & 6.870 & 0.232 \\
10 & 2,047 & 11 & 35.032 & 0.293 \\
12 & 8,191 & 13 & 153.160 & 0.379 \\
14 & 32,767 & 15 & 846.283 & 0.478 \\
\bottomrule
\end{tabular}
\end{center}

% END horn-table.tex


The number of proof nodes in both returned certificates is $d+1$ at every
depth. Thus the search-state reduction is not achieved by silently omitting
necessary proof steps. At depth 14, forward search generated 32,766 rule
instances and demand search generated 14; each yielded the same 15-node chain.
The toy rule order intentionally makes the forward target last at its depth.
A different forward strategy, or the actual \grind\ implementation, can behave
differently.

Correctness testing extends beyond this favorable family. The tests exercise
pure dependency cycles, cycles with a base fact, multiple-premise AND-rules,
wrong goals, invalid proof-parent indices, depth/fact limits, and randomized
finite ground Horn systems compared across both search implementations.
These tests support the implementation's behavior, not a completeness claim
for general theorem proving.

\subsection{Nonlinear cases, including unresolved truths}

The named successful cases are a shifted square, two-variable AM--GM, a cubic
box product, a quartic box product, an equality-ideal identity, a shifted square
with rational margin, two-dimensional Cauchy--Schwarz, a product of three lower
bounds, and three-variable variance. The quartic product under a degree-2 cap
and the Motzkin polynomial remain unresolved. A nearly nonnegative but actually
false shifted square also remains unresolved, as required.

For the generated set, each target is a positive rational sum of six products
of constraints from $x,1-x,y,1-y$, with product degree at most four. This
construction intentionally places the targets inside the chosen certificate
cone. It tests discovery, rational reconstruction, and checking on diverse
coefficient combinations; it does not estimate performance on randomly chosen
mathematical inequalities. The median recorded discovery time was about
21.5 milliseconds and the largest about 103.6 milliseconds for this set.
These are single-run local measurements, not a speed comparison with Lean.

Each of the 24 additional negative controls has the form
$(x-a)^2-\varepsilon$ with $\varepsilon>0$ rational. The benchmark separately
checks its exact negative value at $a$. The cone routine itself returns
\code{unknown}; it does not return a certified refutation. Tests also corrupt
weights, remove sign assumptions, change targets by tiny rational quantities,
and introduce invalid factor indices. The checker must reject those corruptions
even if a numerical optimizer would regard the residual as negligible.

\subsection{Bernstein and invariant evidence}

The Bernstein set contains the worked two-leaf positive polynomial, a box
product, a dyadic zero, a nondyadic zero, a false polynomial, and 40 strictly
positive rational quadratics. The 43 accepted cases have complete checked
subdivision trees. A test corrupts the box-cover geometry to ensure local
positivity does not masquerade as a complete proof.

The invariant set contains nine increments $(n+1)^k$ for $0\leq k\leq8$ and
40 fixed-seed random rational polynomial increments with varied starting
values. Every accepted proposal passes exact base and transition checks.
The median recorded synthesis time was about 4.1 milliseconds. This includes
interpolation and validation but not Lean recursor application, source
elaboration, or kernel checking.

The package exports the nine named cone certificates and all 49 invariant
certificates as JSON. A separate reload-and-validate script accepted all 58
saved records. The additional generated certificates and box trees are
reproducibly rediscovered and checked by the benchmark script; they are not
all serialized as separate proof artifacts.

\subsection{What the Lean export does and does not establish}

The exporter emits nine inequality theorems using explicit nonnegative terms,
rational arithmetic, and polynomial identities. It also emits 49 recurrence
proofs using a generic induction lemma for iteration and the generated exact
base/step expressions. The files are \path{lean/Forge/Certificates.lean} and
\path{lean/Forge/Invariants.lean}.

They contain no deliberate admissions or new axioms. Nevertheless, the source
was not compiled, so it is not a verified Lean deliverable. Elaboration details,
tactic API changes, or dependency integration may require corrections. The
right acceptance criterion is a successful build followed by an axiom and
proof audit, not the plausible appearance of generated source. The included
runner leaves this evidence gap explicit.

\section{A fair evaluation for a real Forge implementation}

\subsection{Benchmark strata and leak prevention}

A subsequent Lean evaluation should use distinct strata: standard \grind\
regressions; nonlinear inequalities; fresh recursive definitions needing
induction; accumulator and fusion lemmas; existential constructions; Boolean
and bit-vector goals; quantified first-order problems; and mixed goals in
which one technique supplies a lemma to another. Keep easy goals, failed goals,
and unsupported goals, not only examples chosen after an extension succeeds.

Partition by theorem family and source module rather than individual nearby
lemmas. Remove the target theorem from any retrieval index. Check for equivalent
lemmas, renamed copies, and specialization leakage. A synthesis benchmark
whose target is already available as a library fact primarily evaluates
retrieval. That can be useful, but it is a different experiment.

Record the complete imported environment, theorem attributes, simplification
sets, seeds, solver versions, and tactic configuration. Reset proof state
between runs; do not let a successful trial add a lemma that helps a later
baseline run.

\subsection{Baselines and ablations}

Compare stock \grind, explicitly configured \grind\ variants, relevant
individual tactics, and a straightforward portfolio. Then compare Forge with
one extension removed at a time: demands, nonlinear cuts, induction
generalization, invariant synthesis, witnesses, or external islands. Compare
the cooperative version with the same engines run independently. Otherwise,
a gain could come entirely from adding a known tactic rather than from the
proposed information exchange.

Use two resource protocols. Under \emph{fixed total resources}, all systems
receive the same wall-clock limit, memory limit, core count, and appropriately
reported heartbeat budget. Under \emph{additive fallback resources}, allow the
unchanged baseline budget plus a declared extension budget, reporting both
budgets and the extra cost. Neither protocol substitutes for the other.

Metrics should include checked solved goals, timeouts, elaboration errors,
unsupported translations, proof bytes, reconstruction time, kernel replay time,
and unexpected axioms. Report paired wins and losses, not only aggregate
success. Repeat timings in randomized order and report confidence intervals
or paired distributions. Include proof checking in the main end-to-end result,
while retaining search-only timing for diagnosis.

The supplied comparison script is much smaller: it tries existing tactics on
nine exported certificate contexts, records compiler diagnostics, and optionally
builds the replay library. It is a smoke-test scaffold, not an implementation
of this full protocol.

\section{Development sequence and acceptance gates}

\textbf{Stage 1: reliable replay.} Compile the generated Lean files, fix any
elaboration issues, and audit their axioms. Implement a typed polynomial
reifier and explicit certificate replay before connecting an external oracle.
Reject malformed certificates and wrong-context assumptions in Lean tests.
A successful stage means that a saved certificate proves the exact intended
Lean theorem, not that discovery is fast.

\textbf{Stage 2: nonlinear integration.} Connect the bounded LP proposal service
to the reifier and checker. Add the goal-sensitive square dictionary, product
cuts, and failure diagnostics. Start with real polynomial inequalities and
explicit hypotheses; postpone complicated casts and rational expressions until
their transport lemmas are tested. Measure against \code{nlinarith} and stock
\grind, including cases they already solve.

\textbf{Stage 3: structural coordination.} Implement the scoped obligation graph,
tabled demands, and one-variable structural induction with changed-argument
generalization. Port additive invariant synthesis and prove its recurrence
transport. Demonstrate mixed examples in which synthesized structure feeds
facts to the algebraic engine. Keep a replay script for every success.

\textbf{Stage 4: richer synthesis and delegation.} Add typed witness grammars,
selected solver islands, and optional broader certificate search. Each backend
must pass a proof-rule coverage matrix and strict/native trust-profile audit.
A missing or incompatible backend should disable only its lane.

\textbf{Stage 5: validated deployment.} Run the controlled repository-scale
benchmarks, tune scheduling on a separate training set, preserve stock behavior
in conservative mode, and publish wins, regressions, unsupported cases, and
proof-checking costs. Do not release an experimental always-on policy merely
because it improves an aggregate score on its tuning examples.

These are dependency-ordered milestones, not a promised schedule. The highest
risk is the typed, scoped, version-compatible bridge into Lean and the quality
of the cooperative search policy, not the small algebraic checker alone.

\section{Conclusion}

The credible path beyond \grind\ is to supply the proof structures and
certificates that local saturation lacks, while retaining its mature core.
Forge makes that path concrete: demand-directed obligations, exact nonlinear
certificates, complete box proofs, justified induction generalization,
polynomial invariants, typed witnesses, and narrowly translated solver islands.
Its organizing rule is that speculative search may be aggressive but every
logical consequence must have a checkable derivation in its actual scope.

The executable work supports four parts of this proposal and exposes meaningful
limitations. It does not establish a finished Lean tactic or measured dominance
over \grind. The next decisive evidence is a typed Lean implementation with
successful kernel replay and a paired, controlled evaluation. The supplied
source, saved certificates, tests, and integration contracts make that a
specific engineering program rather than a vague aspiration.

% END 05-experiments.tex

% BEGIN 06-appendices.tex
\appendix
\section{Certificate formats and independent checking}

\subsection{Sparse polynomials}

A polynomial stores a positive variable count and a sorted list of distinct
exponent vectors with nonzero rational coefficients. Variables are positional.
Exponent vectors have exactly the declared dimension and contain nonnegative
integers. Construction rejects noncanonical encodings instead of silently
merging duplicate entries supplied in a purported certificate.

For example, $x^2-2xy+y^2$ is represented schematically as
\begin{Code}
{
  "nvars": 2,
  "terms": [
    [[0, 2], "1"],
    [[1, 1], "-2"],
    [[2, 0], "1"]
  ]
}
\end{Code}

All fractions in the exported JSON are exact strings such as \code{"1/97"}.
The checker does not parse a decimal into a claim of approximate equality.
An external-facing production parser should additionally bound sizes before
constructing these values.

\subsection{Cone and invariant records}

A cone record contains the target polynomial, ordered lists of nonnegative
and zero hypotheses, and a certificate. Each cone term contains a rational
weight, a square polynomial, and a list of factor indices. An index can occur
more than once. The equality part contains one arbitrary polynomial multiplier
per zero hypothesis. The exact order is part of the certificate's meaning.

An invariant record contains an invariant polynomial, an initial rational state,
and one transition polynomial per state coordinate. The reload checker tests
both base evaluation and exact polynomial substitution. It is independent of
the interpolation and LP discovery functions. It is not independent of the
shared rational-polynomial implementation, a limitation addressed by tests
rather than concealed as formal verification.

The essential cone-checking logic is the following executable pattern:
\begin{Code}
acc = zero_polynomial
for term in certificate.terms:
    reject unless term.weight is rational and nonnegative
    p = term.weight * term.square**2
    for i in term.factors:
        reject unless i is a valid hypothesis index
        p *= nonnegative_hypotheses[i]
    acc += p
for multiplier, equality in paired_ideal_entries:
    acc += multiplier * equality
accept exactly when acc == target
\end{Code}

The actual implementation also checks dimensions, array lengths, canonical
polynomials, and malformed input. Exceptions representing invalid evidence
lead to rejection. Resource-exhaustion behavior is an additional production
requirement, not a theorem proved about the prototype.

\subsection{Horn and box records}

A Horn proof node records its ground conclusion, the rule index or input-fact
marker, a complete ground substitution, and earlier premise-node indices.
The checker verifies the instantiated rule body and head against those exact
nodes. Trimming retains only ancestors of the root, preserving the topological
order. A dependency cycle in the search graph may be legitimate, but the final
proof certificate cannot justify itself through that cycle.

A Bernstein certificate is either a leaf with a box, coordinate degrees, and
all basis coefficients, or a split with an axis, interior rational cut, and
both children. Every basis index must appear exactly once in canonical order.
The checker reconstructs the polynomial at each leaf and checks the complete
box cover recursively.

\section{Reproducing the package}

From the extracted package root, the Python path is:
\begin{Code}
python -m venv .venv
. .venv/bin/activate
python -m pip install -r prototype/requirements.txt
bash scripts/reproduce.sh
\end{Code}

For the exact versions used in this run, substitute
\path{prototype/requirements-tested.txt}. That file is a tested-version
snapshot, not a promise that binary wheels exist for every platform.

Individual actions can be run from \path{prototype/}:
\begin{Code}
python -m pytest -q
python run_benchmarks.py --output ../results
python validate_artifacts.py --results ../results
python export_lean.py --results ../results --output ../lean
\end{Code}

The reproduction script regenerates measurements and certificates; wall-clock
timings naturally vary. It does not download or install Lean. To attempt the
uncompiled replay project on a machine with Lean/Lake and network access:
\begin{Code}
cd lean
lake update
lake exe cache get
lake build
cd ..
python scripts/run_lean_comparisons.py --timeout 30
\end{Code}

The runner preserves compiler output and distinguishes timeout from nonzero
exit. A nonzero compiler exit can be an elaboration or dependency failure,
not necessarily a failed mathematical search. The runtime absence recorded in
the supplied results is intentional and is not counted as either a solve or
a regression.

The article is built with LuaLaTeX. Its main source is
\path{article/forge.tex}, which inputs the numbered chapter files,
\path{horn-table.tex}, and \path{references.tex}. No font files are bundled.
The build helper updates the numerical Horn table from the saved JSON before
running the typesetter.

% END 06-appendices.tex

\clearpage
\addcontentsline{toc}{section}{References and source provenance}
% BEGIN references.tex
\begin{thebibliography}{99}
\small
\bibitem{lean-grind}
Lean developers. \emph{The Lean Language Reference: The grind tactic}.
Current documentation inspected 14 September 2026.
\url{https://lean-lang.org/doc/reference/latest/The--grind--tactic/}.

\bibitem{grind-types}
Lean developers. \emph{Lean.Meta.Tactic.Grind.Types}. Generated API
documentation, inspected 14 September 2026.
\url{https://leanprover-community.github.io/mathlib4_docs/Lean/Meta/Tactic/Grind/Types.html}.

\bibitem{grind-action}
Lean developers. \emph{src/Lean/Meta/Tactic/Grind/Action.lean}, tag v4.34.0.
Inspected source describes the continuation-based control interface,
backtracking, and script reconstruction.
\url{https://github.com/leanprover/lean4/blob/v4.34.0/src/Lean/Meta/Tactic/Grind/Action.lean}.

\bibitem{aesop}
Lean community. \emph{Aesop: White-box automation for Lean 4}.
Repository documentation, inspected 14 September 2026.
\url{https://github.com/leanprover-community/aesop}.

\bibitem{linarith}
Mathlib contributors. \emph{Mathlib.Tactic.Linarith.Frontend}.
Implementation overview and configuration documentation.
\url{https://leanprover-community.github.io/mathlib4_docs/Mathlib/Tactic/Linarith/Frontend.html}.

\bibitem{duper}
Lean community. \emph{Duper}. Repository documentation and feature description,
inspected 14 September 2026.
\url{https://github.com/leanprover-community/duper}.

\bibitem{lean-smt}
A. Mohamed, T. Mascarenhas, H. Khan, H. Barbosa, A. Reynolds, Y. Qian,
C. Tinelli, and C. Barrett.
\emph{lean-smt: An SMT tactic for discharging proof goals in Lean}.
arXiv:2505.15796, version 1, 2025.
\url{https://arxiv.org/html/2505.15796v1}.
Implementation: \url{https://github.com/ufmg-smite/lean-smt}.

\bibitem{interventions}
E. Wang, S. Chess, S. Szeto, and T. Meek.
\emph{Learned Interventions Inside Lean 4's grind}.
arXiv:2607.22972, version 1, 2026.
\url{https://arxiv.org/html/2607.22972v1}.
Cited as a research report; its experiments are not reproduced here.

\bibitem{lean-release}
Lean developers. \emph{Lean v4.34.0 release}.
Published 14 September 2026; release metadata inspected through the GitHub API.
\url{https://github.com/leanprover/lean4/releases/tag/v4.34.0}.

\bibitem{mathlib-revision}
Mathlib contributors. Repository revision
\texttt{1cf325a0cf67aca2b04d76b5380ff6a9e410aefa} and its
\texttt{lean-toolchain} file.
\url{https://github.com/leanprover-community/mathlib4/tree/1cf325a0cf67aca2b04d76b5380ff6a9e410aefa}.

\bibitem{scipy-lp}
SciPy developers. \emph{scipy.optimize.linprog}.
Official API reference. The local experiment used SciPy 1.17.0 with the
\texttt{highs} method, not necessarily the version currently displayed by the site.
\url{https://docs.scipy.org/doc/scipy/reference/generated/scipy.optimize.linprog.html}.

\bibitem{lean-tactics}
Lean developers. \emph{The Lean Language Reference: Tactic Reference}.
Sections on \texttt{native\_decide}, \texttt{decide\_cbv}, and \texttt{omega},
inspected 14 September 2026.
\url{https://lean-lang.org/doc/reference/latest/Tactic-Proofs/Tactic-Reference/}.

\bibitem{ring}
Mathlib contributors. \emph{Mathlib.Tactic.Ring.Basic}.
\url{https://leanprover-community.github.io/mathlib4_docs/Mathlib/Tactic/Ring/Basic.html}.

\bibitem{positivity}
Mathlib contributors. \emph{Mathlib.Tactic.Positivity.Basic}.
\url{https://leanprover-community.github.io/mathlib4_docs/Mathlib/Tactic/Positivity/Basic.html}.
\end{thebibliography}

\medskip
\textbf{Provenance note.} External sources establish the baseline, available
interfaces, related work, and library behavior. The architecture, derivations,
Python code, tests, and reported measurements in this package were developed
for this task. Research inspiration is distinguished from executed integration;
no third-party research paper, repository archive, or font file is redistributed.

% END references.tex

\end{document}
